\BourbakiExercisesSection

\begin{enumerate}
\def\labelenumi{\arabic{enumi}.}
\item
  Which are the field structures amongst the ring structures determined in § 8, Exercise 1?
\item
  A finite ring in which every non-zero element is cancellable is a field (cf.~§ 2, Exercise 6).
\item
  Let G be a commutative group with operators which is simple (§ 4, no. 4, Definition 7). Show that the endomorphism ring A of G is a field.
\item
  Let K be a field. Show that there exists a smallest subfield F of K and that F is canonically isomorphic either to the field Q or to the field Z/pZ for some prime number p.~In the first case (resp. second case) K is said to be of characteristic 0 (resp. p).
\item
  Let K be a commutative field of characteristic \(\neq2\) ; let G be a subgroup of the additive group of K such that, if H denotes the set consisting of 0 and the inverses of the elements \(\neq0\) of G, H is also a subgroup of the additive group of K. Show that there exists an element \(a\in K\) and a subfield \(K'\) of K such that \(G=aK'\) (establish first that, if x and y are elements of G such that \(y\neq0\) , then \(x^{2}/y\in G\) ; deduce that, if x,y,z are elements of G such that \(z\neq0\) , then xy/z\ensuremath{\in}G).
\item
  Let K be a commutative field of characteristic \(\neq2\) ; let f be a mapping of K into K such that \(f(x+y)=f(x)+f(y)\) for all x and y and for all \(x\neq0\) \(f(x)f(1/x)=1\) . Show that f is an isomorphism of K onto a subfield of K (prove that \(f(x^{2})=(f(x))^{2}\) ).
\item
  Let A be a commutative ring.
\end{enumerate}

\begin{enumerate}
\def\labelenumi{(\alph{enumi})}
\item
  Every prime ideal is irreducible (§ 8, Exercise 11).
\item
  The set of prime ideals is inductive for the relations \(\subset\) and \(\supset\) .
\item
  Let \(\mathfrak{a}\) be an ideal of \(\mathbf{A}\) distinct from \(\mathbf{A}\) ; let \(\mathfrak{b}\) be the set of \(x \in \mathbf{A}\) such that there exists an integer \(n \geqslant 0\) with \(x^n \in \mathfrak{a}\) ( \(n\) depends on \(x\) ). Show that \(\mathfrak{b}\) is an ideal and is equal to the intersection of the prime ideals of \(\mathbf{A}\) containing \(\mathfrak{a}\) .
\end{enumerate}

¶ 8. A ring A is called a Boolean ring if each of its elements is idempotent (in other words, if \(x^2 = x\) for all \(x \in \mathbf{A}\) ).

\begin{enumerate}
\def\labelenumi{(\alph{enumi})}
\item
  In the set \(\mathfrak{P}(\mathbf{E})\) of subsets of a set \(\mathbf{E}\) , show that a Boolean ring structure is defined by setting \(\mathbf{AB} = \mathbf{A} \cap \mathbf{B}\) and \(\mathbf{A} + \mathbf{B} = (\mathbf{A} \cap \mathbb{C}\mathbf{B}) \cup (\mathbf{B} \cap \mathbb{C}\mathbf{A})\) . This ring is isomorphic to the ring \(\mathbf{K}^{\mathrm{E}}\) of mappings of \(\mathbf{E}\) into the ring \(\mathbf{K} = \mathbf{Z}/(2)\) of integers modulo 2 (consider for each subset \(\mathbf{X} \subset \mathbf{E}\) its ``characteristic function'' \(\phi_{\mathbf{X}}\) such that \(\phi_{\mathbf{X}}(x) = 1\) if \(x \in \mathbf{X}\) , \(\phi_{\mathbf{X}}(x) = 0\) if \(x \notin \mathbf{X}\) ).
\item
  Every Boolean ring \(\mathbf{A}\) is commutative and such that \(x + x = 0\) for \(x \in \mathbf{A}\) (write \(x + x\) as an idempotent and then \(x + y\) as an idempotent).
\item
  If a Boolean ring \(\mathbf{A}\) contains no divisor of 0, it is reduced to 0 or is isomorphic to \(\mathbf{Z} / (2)\) (if \(x\) and \(y\) are any two elements of \(\mathbf{A}\) , show that \(xy(x + y) = 0\) ). Deduce that in a Boolean ring every prime ideal is maximal.
\item
  In a Boolean ring A, every ideal \(a \neq A\) is the intersection of the prime ideals containing a (apply Exercise 7(c)). Deduce that every irreducible ideal is maximal (in other words, that the notions of irreducible ideal, prime ideal and maximal ideal coincide in a Boolean ring).
\item
  Show that every Boolean ring is isomorphic to a subring of a product ring \(\mathbf{K}^{\mathbb{E}}\) , where \(\mathbf{K} = \mathbf{Z}/(2)\) (use \((d)\) ).
\item
  If \(p_{i}\) ( \(1 \leqslant i \leqslant n\) ) are n distinct maximal ideals in a Boolean ring A and \(a = \bigcap_{1 \leqslant i \leqslant n} p_{i}\) , show that the quotient ring A/a is isomorphic to the product ring \(K^{n}\) . Deduce that every finite Boolean ring is of the form \(K^{n}\) .
\item
  In a Boolean ring A, the relation \(xy = x\) is an order relation; if it is denoted by \(x \leqslant y\) , show that with this relation A is a distributive lattice (Set Theory, III, § 1, Exercise 16), has a least element \(\alpha\) and that, for every ordered pair \((x, y)\) of elements such that \(x \leqslant y\) , there exists an element \(d(x, y)\) such that \(\inf(x, d(x, y)) = \alpha\) , \(\sup(x, d(x, y)) = y\) . Conversely, if an ordered set A has these three properties, show that the laws of composition \(xy = \inf(x, y)\) and \(x + y = d(\sup(x, y), \inf(x, y))\) define a Boolean ring structure on A.
\end{enumerate}

\begin{enumerate}
\def\labelenumi{\arabic{enumi}.}
\setcounter{enumi}{8}
\tightlist
\item
  Let \(\mathbf{A}\) be a ring such that \(x^3 = x\) for all \(x \in \mathbf{A}\) . It is proposed to prove that \(\mathbf{A}\) is commutative.
\end{enumerate}

\begin{enumerate}
\def\labelenumi{(\alph{enumi})}
\item
  Show that \(6\mathbf{A} = \{0\}\) and that \(2\mathbf{A}\) and \(3\mathbf{A}\) are two-sided ideals such that \(2\mathbf{A} + 3\mathbf{A} = \mathbf{A}\) and \(2\mathbf{A} \cap 3\mathbf{A} = \{0\}\) . Deduce that it can be assumed that either \(2\mathbf{A} = \{0\}\) or \(3\mathbf{A} = \{0\}\) .
\item
  If 2A = \{0\}, calculate \((1 + x)^{3}\) , deduce that \(x^{2} = x\) for all \(x \in A\) and conclude by means of Exercise 8.
\item
  If \(3\mathbf{A} = \{0\}\) , calculate \((x + y)^3\) and \((x - y)^3\) , show that \(x^2y + xyx + yx^2 = 0\) and left multiply by \(x\) ; deduce that \(xy - yx = 0\) .
\item
  Let A be a ring such that \(3A = \{0\}\) and \(x^{3} = x\) for all \(x \in A\) ; define a set I and an injective homomorphism of A into \((\mathbf{Z}/3\mathbf{Z})^{\mathrm{I}}\) (same method as in Exercise 8).
\end{enumerate}

\begin{enumerate}
\def\labelenumi{\arabic{enumi}.}
\setcounter{enumi}{9}
\tightlist
\item
  Let \(\mathbf{A}\) be a commutative ring and \(\mathfrak{U}\) the intersection of its maximal ideals.
\end{enumerate}

\begin{enumerate}
\def\labelenumi{(\alph{enumi})}
\item
  Show that the canonical homomorphism \(\mathbf{A}^* \to (\mathbf{A} / \mathfrak{U})^*\) is surjective and that its kernel is \(1 + \mathfrak{U}\) .
\item
  Suppose that the set M of maximal ideals of A is finite and that A* is finite. Deduce that A is finite (apply Proposition 8 of §8, no. 10 to A/U).
\item
  Deduce that the set of prime numbers is infinite (note that \(\mathbf{Z}^{*} = \{1, -1\}\) ).
\item
  Suppose that A* is finite and that A/m is finite for every maximal ideal m. Can it be deduced that A is finite, that A is countable?
\end{enumerate}

\begin{enumerate}
\def\labelenumi{\arabic{enumi}.}
\setcounter{enumi}{10}
\tightlist
\item
  Let P be the set of prime numbers and A the product ring of the fields \(\mathbf{Z} / p\mathbf{Z}\) , \(p \in \mathbb{P}\) . Let \(a\) be the subset of A consisting of the elements \((a_p)_{p \in \mathbb{P}}\) such that \(a_p \neq 0\) only for a finite number of indices \(p\) .
\end{enumerate}

\begin{enumerate}
\def\labelenumi{(\alph{enumi})}
\item
  a is an ideal of A.
\item
  Let \(\mathbf{B} = \mathbf{A} / a\) . For every integer \(n > 0\) and every \(b \neq 0\) in \(\mathbf{B}\) , there exists one and only one element \(b'\) of \(\mathbf{B}\) such that \(nb' = b\) (note that if \(p \in \mathbb{P}\) does not divide \(n\) , multiplication by \(n\) in \(\mathbf{Z} / p\mathbf{Z}\) is bijective). Deduce that \(\mathbf{B}\) contains a subfield isomorphic to \(\mathbf{Q}\) .
\end{enumerate}

\begin{enumerate}
\def\labelenumi{\arabic{enumi}.}
\setcounter{enumi}{11}
\item
  Let A be a commutative ring and \(a, b \in A\) . Show that the canonical image of ab in \(\mathrm{A}/(a - a^{2}b)\) is an idempotent. Give an example where this idempotent is distinct from 0 and 1.
\item
  Determine the endomorphisms of the ring Z and the ring \(Z \times Z\) . More generally, if I and J are finite sets, determine the homomorphisms of the ring \(Z^{1}\) into the ring \(Z^{J}\) .
\item
  Let A and B be two rings and f and g two homomorphisms of A into B. Is the set of \(x \in A\) such that \(f(x) = g(x)\) a subring of A, an ideal of A?
\end{enumerate}

¶ 15. Let A be a non-commutative ring with no divisors of 0. A is said to admit a field of left fractions if it is isomorphic to a subring B of a field K such that every element of K is of the form \(x^{-1}y\) , where \(x \in \mathbf{B}, y \in \mathbf{B}\) .

\begin{enumerate}
\def\labelenumi{(\alph{enumi})}
\tightlist
\item
  Let A' be the set of elements \(\neq0\) in A. For A to admit a field of left fractions, it is necessary that the following condition be fulfilled:
\end{enumerate}

\begin{enumerate}
\def\labelenumi{(\Alph{enumi})}
\setcounter{enumi}{6}
\tightlist
\item
  for all \(x \in \mathbf{A}\) , \(x' \in \mathbf{A}'\) , there exist \(u \in \mathbf{A}'\) and \(v \in \mathbf{A}\) such that \(ux = vx'\) .
\end{enumerate}

\begin{enumerate}
\def\labelenumi{(\alph{enumi})}
\setcounter{enumi}{1}
\tightlist
\item
  Suppose conversely that condition (G) is fulfilled. Show that in the set A × A' the relation R between \((x, x')\) and \((y, y')\) which states ``for every ordered pair \((u, v)\) of elements ≠0 such that \(ux' = vy'\) , \(ux = vy\) is an equivalence relation.
\end{enumerate}

Let \((x, x')\) , \((y, y')\) be two elements of \(A \times A'\) , \(\xi\) and \(\eta\) their respective classes (mod. R). For every ordered pair \((u, u') \in A \times A'\) such that \(u'x = uy'\) , show that the class (mod. R) of \((uy, u'y')\) depends only on the classes \(\xi\) and \(\eta\) ; if it is denoted by \(\xi\eta\) , a law of composition is defined on the set \(K = (A \times A')/R\) . If \(K'\) is the set of elements of K distinct from the class 0 of elements \((0, x')\) of \(A \times A'\) , \(K'\) , with the law induced by the above law, is a group.

For every element \(x \in \mathbf{A}\) , the elements \((x'x, x')\) , where \(x'\) runs through \(\mathbf{A}'\) , define an isomorphism of \(\mathbf{A}\) (with multiplication alone) onto a subring of \(\mathbf{K}\) . Identifying \(\mathbf{A}\) with its image under this isomorphism, the class (mod. R) of an ordered pair \((x, x') \in \mathbf{A} \times \mathbf{A}'\) is identified with the element \(x'^{-1}x\) .

This being so, if \(\xi = x'^{-1}x\) is an element of K and 1 denotes the unit element of \(\mathbf{K}'\) , then \(\xi + 1\) denotes the element \(x'^{-1}(x + x')\) , which does not depend on the representation of \(\xi\) in the form \(x'^{-1}x\) . We then write \(\xi + 0 = \xi\) and, for \(\eta \neq 0\) , \(\xi + \eta = \eta(\eta^{-1}\xi + 1)\) . Show that the addition and multiplication thus defined on K determine on that set a field structure which extends the ring structure on A; in other words, condition (G) is sufficient for A to admit a field of left fractions.

\begin{enumerate}
\def\labelenumi{\arabic{enumi}.}
\setcounter{enumi}{15}
\item
  Let A be a ring in which every non-zero element is cancellable. For A to admit a field of left fractions (Exercise 15), it is necessary and sufficient that in A the intersection of two left ideals distinct from 0 never reduce to 0.
\item
  Let \((\mathbf{K}_i)_{i\in I}\) be a family of fields. For every subset \(J\) of \(I\) , let \(e_J\) denote the element of the product \(\mathbf{K} = \prod_{i\in I}\mathbf{K}_i\) whose \(i\) -th component is equal to 0 if \(i\in J\) and to 1 if \(i\in I - J\) .
\end{enumerate}

\begin{enumerate}
\def\labelenumi{(\alph{enumi})}
\item
  Let \(x = (x_{i})\) be an element of \(\mathbf{K}\) and \(J(x)\) the set of elements \(i \in I\) such that \(x_{i} = 0\) . Show that there exists an invertible element \(u\) in \(\mathbf{K}\) such that \(x = ue_{J(x)} = e_{J(x)}u\) .
\item
  Let \(\mathfrak{a}\) be a left (resp. right) ideal of \(\mathbf{K}\) distinct from \(\mathbf{K}\) . Let \(\mathfrak{F}_{\mathfrak{a}}\) be the set of subsets \(\mathbf{J}\) of \(\mathbf{I}\) such that \(e_{J} \in \mathfrak{a}\) . Show that \(\mathfrak{F}_{\mathfrak{a}}\) has the following properties:
\end{enumerate}

\((b_{1})\) If \(\mathbf{J} \in \mathfrak{F}_{\mathfrak{a}}\) and \(\mathbf{J}' \supset \mathbf{J}\) , then \(\mathbf{J}' \in \mathfrak{F}_{\mathfrak{a}}\) .

\((b_{2})\) If \(\mathrm{J}_1, \mathrm{J}_2 \in \mathfrak{F}_{\mathfrak{a}}\) , then \(\mathrm{J}_1 \cap \mathrm{J}_2 \in \mathfrak{F}_{\mathfrak{a}}\) .

\[
\left(b _ {3}\right) \varnothing \notin \mathfrak {F} _ {\mathfrak {a}}.
\]

Show that \(x \in \mathfrak{a}\) is equivalent to \(J(x) \in \mathfrak{F}_{\mathfrak{a}}\) (use (a)). Deduce that \(\mathfrak{a}\) is a two-sided ideal.

\begin{enumerate}
\def\labelenumi{(\alph{enumi})}
\setcounter{enumi}{2}
\item
  A set of subsets of I is called a filter if it satisfies properties \((b_{1})\) , \((b_{2})\) , \((b_{3})\) above (*cf.~General Topology, I, §6*). Show that the mapping \(\mathfrak{a} \mapsto \mathfrak{F}_{\mathfrak{a}}\) is a strictly increasing bijection of the set of ideals of K distinct from K onto the set of filters of I. *Show that a is maximal if and only if \(\mathfrak{F}_{\mathfrak{a}}\) is an ultrafilter.*
\item
  *Let \(\mathfrak{F}\) be a non-trivial ultrafilter on the set \(P\) of prime numbers, let \(a\) be the corresponding ideal of the ring \(\mathbf{K} = \prod_{p \in P} \mathbf{Z}/p\mathbf{Z}\) and let \(k = K/a\) . Show that \(k\) is a field of characteristic \(0\)*
\end{enumerate}

¶ 18. (a) Let K be a field and a, b two non-permutable elements of K. Show that

\[
\begin{array}{l} (1) a = (b - (a - 1) ^ {- 1} b (a - 1)) (a ^ {- 1} b a - (a - 1) ^ {- 1} b (a - 1)) ^ {- 1} \\ (2) a = (1 - (a - 1) ^ {- 1} b ^ {- 1} (a - 1) b) (a ^ {- 1} b ^ {- 1} a b - (a - 1) ^ {- 1} b ^ {- 1} (a - 1) b) ^ {- 1}. \end{array}
\]

\begin{enumerate}
\def\labelenumi{(\alph{enumi})}
\setcounter{enumi}{1}
\item
  Let K be a non-commutative field and x an element not belonging to the centre of K. Show that K is generated by the set of conjugates \(axa^{-1}\) of x. (Let \(K_{1}\) be the subfield of K generated by the set of conjugates of x. Deduce from (a) that, if \(K_{1} \neq K\) and \(a \in K_{1}\) and \(b \notin K_{1}\) , then necessarily ab = ba. Deduce a contradiction by considering in \(K_{1}\) two non-permutable elements a, \(a'\) and an element \(b \notin K_{1}\) and noting that \(ba \notin K_{1}\) .)
\item
  Deduce from (b) that if Z is the centre of K and H is a subfield of K such that \(Z \subset H \subset K\) and \(aHa^{-1} = H\) for all \(a \neq 0\) in K, then H = Z or H = K (Cartan-Brauer-Hua Theorem).
\item
  Deduce from (a) that in a non-commutative field K, the set of commutators \(aba^{-1}b^{-1}\) of elements \(\neq0\) in K generates the field K.
\end{enumerate}

\begin{enumerate}
\def\labelenumi{\arabic{enumi}.}
\setcounter{enumi}{18}
\tightlist
\item
  Let A be a non-zero pseudo-ring such that, for all \(a \neq 0\) in A and all \(b \in \mathbf{A}\) , the equation \(ax + ya = b\) has solutions consisting of ordered pairs \((x, y)\) of elements of A; in other words, \(a\mathbf{A} + \mathbf{A}a = \mathbf{A}\) for \(a \neq 0\) in A.
\end{enumerate}

\begin{enumerate}
\def\labelenumi{(\alph{enumi})}
\tightlist
\item
  Show that A is the only non-zero two-sided ideal in A.
\end{enumerate}

\begin{enumerate}
\def\labelenumi{(\alph{enumi})}
\setcounter{enumi}{1}
\item
  Show that the relation \(a \neq 0\) implies \(a^2 \neq 0\) (note that if \(a^2 = 0\) , it would follow that \(aAa = \{0\}\) ).
\item
  Show that if \(ab = 0\) and \(a \neq 0\) (resp. \(b \neq 0\) ) then \(b = 0\) (resp. \(a = 0\) ). (Note that \((ba)^2 = 0\) and deduce that \(bAb = \{0\}\) .)
\end{enumerate}

\begin{enumerate}
\def\labelenumi{\arabic{enumi}.}
\setcounter{enumi}{19}
\item
  Let A be a non-zero pseudo-ring such that there exists \(a \in A\) for which, for all \(b \in A\) , one of the equations ax = b, xa = b admits a solution \(x \in A\) . Show that if A has no divisor of 0 other than 0, A admits a unit element.
\item
  Let A be a non-zero pseudo-ring in which, for all \(a \neq 0\) and all \(b \in A\) , one of the equations ax = b, xa = b admits a solution \(x \in A\) . Show that A is a field (use Exercises 19 and 20).
\end{enumerate}

